% Context from chapters-tex/convex-analysis.tex; for mathematical reference, not compilation.
tion of half-spaces, it is a closed convex set.
                                                      
 
Thus $\partial f : \Hh \mapsto 2^{\Hh^*}$ is a set-valued operator, also written $\partial f : \Hh \hookrightarrow \Hh^*$.

\begin{rem}[Maximally monotone operator]
The subdifferential $\partial f$ is monotone: setting $U=\partial f$ gives
\eq{
	\foralls (u,v) \in U(x) \times U(y), \quad
		\dotp{y-x}{v-u} \geq 0. 
}
For single-valued maps on the real line, monotonicity means being nondecreasing.
 
Subdifferentials can also be shown to be maximally monotone, meaning that their graphs are not properly contained in the graph of another monotone operator.
 
Conversely, a monotone map need not be a subdifferential; an example is $(x,y) \mapsto (-y,x)$.
  
Many results extend from subdifferentials to arbitrary maximally monotone operators. We restrict attention to subdifferentials here.
\end{rem}

For the absolute value $f(x)=|x|$, write $\partial f(x) = \partial |\cdot|(x)$. Its subdifferential is
\eq{
	\partial |\cdot|(x) = 
	\choice{
		\{-1\} \qifq x < 0, \\
		\{1\} \qifq x > 0, \\
		{[-1,1]} \qifq x=0.
	}
}

\begin{figure}[tbp]
\centering
\BookDrawing{tikz/convex-analysis-subdiff-l1}
\caption{Subdifferential of the absolute value and a piecewise affine convex function.}
\label{fig-subdiff-l1}
\end{figure}







   
\paragraph{First-order conditions.}

The subdifferential gives a necessary and sufficient optimality condition.

\begin{prop}
	$x^\star$ is a minimizer of $f$ if and only if $0 \in \partial f(x^\star)$.
\end{prop}
\begin{proof}
	One has
	\eq{
		x^\star \in \argmin f
		\quad\Leftrightarrow\quad
		\pa{ \forall y, f(x^\star) \leq f(y) + \dotp{0}{x^\star-y} }
		\quad\Leftrightarrow\quad
		0 \in \partial f(x^\star).
	}
\end{proof}

   
\paragraph{Subdifferential calculus.}

Calculus rules simplify subdifferential computations. For a separable function $f(x_1,\ldots,x_K)=\sum_{k=1}^K f_k(x_k)$, the subdifferential is the Cartesian product
\eq{
	\partial f(x_1,\ldots,x_K) = \partial f_1(x_1) \times \ldots \times \partial f_K(x_K).
}
Applying this rule to the $\ell^1$ norm $\norm{x}_1=\sum_{k=1}^N |x_k|$ gives
\eq{
	\partial \norm{\cdot}_1(x) = \prod_{k=1}^N \partial |\cdot|(x_k)
}
a hyperrectangle. With $I = \supp(x)$, the condition $u \in \partial \norm{\cdot}_1(x)$ is equivalent to
\eq{
	u_I = \sign(x_I)
	\qandq
	\norm{u_{I^c}}_\infty \leq 1.
}

A sum rule requires a domain qualification. If one function is finite and continuous at a point in the domain of the other, then for every $x\in\dom(f)\cap\dom(g)$,
\eq{
	\partial (f+g)(x) = \partial f(x) \oplus \partial g(x) = \enscond{u+v}{(u,v) \in \partial f(x) \times \partial g(x)}
}
where $\oplus$ denotes the Minkowski sum. In particular, if $f$ is differentiable at $x$, then
\eq{
	\partial (f+g)(x) = \nabla f(x) + \partial g(x) = \enscond{\nabla f(x) + v }{ v \in \partial g(x) }.
}
Positive scaling gives
\eq{
	\foralls \la >0, \quad
	\partial (\la f)(x) = \la (\partial f(x)). 
}

General compositions need not preserve convexity, so no unrestricted chain rule is available.
 
Composition with a linear map does preserve convexity. For $A \in \RR^{P \times N}$ and $f \in \Ga_0(\RR^P)$, assume $\Im(A)\cap\relint(\dom(f))\neq\emptyset$. Then $f \circ A \in \Ga_0(\RR^N)$ and
\eq{
	\partial (f \circ A)(x) = A^* (\partial f)(Ax) \eqdef \enscond{A^* u}{ u \in \partial f(Ax) }. 
}

\begin{rem}[Application to the Lasso]
For $F(x)=\frac12\norm{Ax-y}^2+\lambda\norm{x}_1$ with $\lambda>0$, the sum rule gives
\eq{0\in A^*(Ax-y)+\lambda\partial\norm{\cdot}_1(x).}
Writing $r=y-Ax$, optimality is equivalent to $(A^*r)_i=\lambda\sign(x_i)$ on $\supp(x)$ and $|(A^*r)_i|\leq\lambda$ elsewhere. These conditions are useful both for analysis and for checking a numerical solution.
\end{rem}

   
\paragraph{Normal cone.}


For a nonempty closed convex set $\Cc$, the subdifferential of its indicator is the normal cone:
\eq{
	\foralls x \in \Cc, \quad \partial \iota_\Cc(x) = \Nn_\Cc(x) \eqdef \enscond{ v }{ \foralls z \in \Cc, \dotp{z-x}{v} \leq 0 }.
}
Outside the set, $x \notin \Cc$ implies $\partial  \iota_\Cc(x) = \emptyset$.
 
For an affine space $\Cc = a+\Vv$ with direction space $\Vv \subset \Hh$, the normal cone at each of its points is $\Nn_\Cc(x)=\Vv^\bot$, the orthogonal complement of $\Vv$. More generally, at an interior point $x \in \inter(\Cc)$ of $\Cc$, we have $\Nn_\Cc(x)=\{0\}$. At a boundary point, the normal cone describes the supporting directions of $\Cc$ at $x$.






\begin{figure}[tbp]
\centering
\BookDrawing{tikz/convex-analysis-normal-cones}
\caption{Normal cones}
\label{fig-review-convex-analysis-normal-cones}
\end{figure}
The normal cone expresses first-order optimality for the constrained problem
\eq{
	\umin{x \in \Cc} f(x)
}
If $f$ is finite and continuous at a point of $\Cc$, the sum rule gives the optimality condition
\eq{
	0 \in \partial f(x) + \partial \iota_{\Cc}(x)
	\quad\Leftrightarrow\quad
	\exists \xi \in \partial f(x), - \xi \in \Nn_\Cc(x)
	\quad\Leftrightarrow\quad
	\partial f(x) \cap (-\Nn_\Cc(x)) \neq \emptyset.
}
If $f$ is differentiable, it reads $-\nabla f(x) \in \Nn_\Cc(x)$.



                                                                                  
                                                                                  
                                                                                  
\section{Legendre--Fenchel Transform}

The Legendre--Fenchel transform provides a dual representation of a convex function. It makes many Lagrange-duality calculations systematic by expressing inner minimizations through conjugates.

                                                                                  
\subsection{Legendre Transform}

For $f \in \Ga_0(\Hh)$, we define its Legendre--Fenchel transform, or convex conjugate, by
\eql{\label{eq-fenchel-transf}
	f^*(u) \eqdef \usup{x} \dotp{x}{u} - f(x).
}
As a supremum of continuous affine functions, $f^*$ is convex and lower semicontinuous; in fact, $f^\star \in \Ga_0(\Hh^*)$. The following biconjugacy theorem recovers the original function.

\begin{thm}
One has
\eq{
	\foralls f \in \Ga_0(\Hh), \quad (f^{*})^* = f. 
}
\end{thm}

Even when $f$ is not convex, $f^*$ is convex. If $f$ admits an affine minorant, $f^{**}$ is its lower semicontinuous convex envelope: the largest lower semicontinuous convex function below $f$.


The subdifferentials of $f$ and $f^*$ satisfy the following inverse relation.

\begin{prop}
We have $\partial f^* = (\partial f)^{-1}$, where the inverse of a set-valued map is defined in~\eqref{eq-inv-setvalued}. The Fenchel--Young inequality states that
\eq{
	\foralls (